\documentclass[tate,paper=a5,fontsize=10pt]{jlreq}
\begin{document}
\section{駅まで}
朝から細かい雨が降っていた。傘を持たずに家を出たのは失敗だったと、駅までの坂道を歩きながら何度も思った。商店街の軒先を伝って歩けば少しは濡れずに済むが、それでも肩のあたりはすぐに冷たくなった。

改札の前で友人が待っていた。彼は大きな黒い傘を畳みながら、遅いぞと笑った。約束の時刻にはまだ五分あったが、彼はいつも十分前には着いている。

\section{車内}
電車は思ったより空いていた。窓際の席に座ると、外の景色が雨粒で滲んで見えた。友人は鞄から文庫本を取り出し、しおりを挟んだ頁を開いた。

例の場所というのが何を指すのか、私にはすぐにわかった。学生の頃によく通った川沿いの喫茶店である。店主は無口な人で、注文を聞く以外にはほとんど口をきかなかった。

木曜日、私は約束の時刻より少し早く店に入った。店内は昔と変わらず薄暗く、奥の席には誰もいなかった。珈琲を頼み、窓の外の川を眺めていると、扉の鈴が鳴った。
\end{document}
